\documentclass[11.5pt,a4paper]{book}

% Inizia il preambolo.

\setlength{\textheight}     {20.0  cm}
\setlength{\textwidth}      {15.0  cm}
\setlength{\baselineskip}   { 1.0  cm}
\setlength{\oddsidemargin}  { 0.4  cm}
\setlength{\evensidemargin} { 0.4  cm}
\setlength{\topmargin}      { 0.4  cm}
\setlength{\parskip}        { 4.0  pt}
\setlength{\arrayrulewidth} { 0.5  pt}
\setlength{\doublerulesep}  { 0.0  pt}
\setlength{\footskip}       {60.0  pt}
\renewcommand{\baselinestretch}{1.25}
\renewcommand{\footnotesize} 
\renewcommand{\bfseries}
\usepackage[dvips]{graphicx}
\usepackage[italian]{babel}
\usepackage{fancyhdr}
\fancyhead[R]{}
\usepackage[usenames]{color}
\usepackage[hang,small,bf]{caption}
\usepackage{amsmath}
\usepackage{amssymb}
\usepackage{braket}
\usepackage{subfig}
\usepackage{multirow}
\usepackage{wrapfig}
\usepackage{siunitx}
\usepackage{textcomp}
\usepackage{longtable}
\usepackage{booktabs}
\usepackage{ucs}
\usepackage[utf8x]{inputenc}
\usepackage{hyperref}
\hypersetup{pdfauthor={Antonio Salvucci},%
  pdftitle={NmapSi4 - Documentation}}
\usepackage{listings}
\lstset{language=Python,%
  basicstyle=\small}
% Fine del preambolo.

\begin{document}

\pagestyle{empty}
%%%%%%%%%%%%%%%%%%%%%%%%%%%%%%%%%%%%%%%%%%%%%%%%%%%%%%%%%%%%%%%%%%%%%%%%%%%%%%
%                       NmapSI4 Latex Documentation                          %
%              copyright: Antonio Salvucci <s4lb04@gmail.com>                %
%%%%%%%%%%%%%%%%%%%%%%%%%%%%%%%%%%%%%%%%%%%%%%%%%%%%%%%%%%%%%%%%%%%%%%%%%%%%%%
{
  \centering
  \includegraphics[width=0.7\textwidth]{nmapsi4}\\
  \LARGE{\bf\boldmath Manuale d'uso}
  
  
  %\Large Author: Antonio Salvucci
}
\pagebreak 

%$\ $

%\pagebreak

\pagenumbering{roman}
\pagestyle{plain}
\unboldmath
\tableofcontents
\include{introd}
\pagenumbering{arabic}
\pagestyle{fancy}
\include{lhc-atlas}
\include{zprime-theory}
\include{alignments}
\include{misalign-effects}
\include{discovery}
\include{concl}
\include{appendix}
\include{thanks}
\include{biblio}
%
\end{document}
